\section{Conclusion}

% \TODO{DH: we can skip the summary here and just start with We evaluate XX and introduce YYY... so no history recap here to save space} Self-intersections in meshes have been studied for decades, and a large number of solutions and repair tools have been proposed. If one's sole objective is to eliminate intersections, in other words, reduce the count to zero, many existing methods--particularly those that subdivide problematic facets--would perform adequately. However, when mesh quality is equally crucial, these standard approaches may fall short. In such contexts, improving mesh quality becomes necessary, potentially through remeshing the affected regions where self-intersections were observed. Moreover, there are scenarios in which simple mesh repair is not sufficient. When it is important to distinguish between different submeshes during intersection resolution, detaching or controlled cutting may be required to preserve their structural meaning.
% The solutions implemented in this study are preliminary prototype designed to explore various methods to address these challenges. Future work will include enhancing its efficiency and robustness of the algorithm, validating it across more datasets, re-implementing it in C++ to improve performance. Ultimately, further research and mesh libraries development should also aim to provide ways that consider mesh quality while resolving self-intersections.

We evaluated several existing methods to merge neighboring meshes that intersect and introduced a new approach. Even when the repaired mesh was geometrically correct, the results were structurally and qualitatively different, hence the optimal method depends on the application, and quality metrics need to be carefully evaluated. Our method resulted in geometrically correct meshes of high quality. Furthermore, we proposed the first successful strategies for detaching intersecting meshes in cases where it is important to distinguish between different substructures. 
The solutions implemented in this study are preliminary prototypes. Future work will include enhancing their quality, efficiency and robustness, validating across more datasets. Further research and mesh library development should also aim to provide ways that improve mesh quality while resolving self-intersections.


